\documentclass[9pt,a4paper]{extarticle}
\usepackage{parskip}

% =====================
% Codificação, idioma e layout
% =====================
\usepackage[T1]{fontenc}
\usepackage[utf8]{inputenc}
\usepackage[brazil]{babel}
\usepackage{lmodern}
\usepackage{geometry}
\geometry{margin=2cm}
\usepackage{microtype}

% =====================
% Matemática e símbolos
% =====================
\usepackage{amsmath, amssymb, amsthm}
\usepackage{mathtools}
\usepackage{bm}

% =====================
% Figuras, tabelas, links
% =====================
\usepackage{graphicx}
\usepackage{booktabs}
\usepackage{hyperref}
\hypersetup{colorlinks=true,linkcolor=blue,citecolor=blue,urlcolor=blue}

% =====================
% Código-fonte
% =====================
\usepackage{listings}
\usepackage{xcolor}
\definecolor{codebg}{rgb}{0.96,0.96,0.96}
\definecolor{codekw}{RGB}{33,96,171}
\definecolor{codecmt}{RGB}{120,120,120}
\definecolor{codestr}{RGB}{163,21,21}
\lstdefinestyle{pystyle}{
  language=Python,
  backgroundcolor=\color{codebg},
  basicstyle=\ttfamily\footnotesize,
  keywordstyle=\color{codekw}\bfseries,
  commentstyle=\color{codecmt}\itshape,
  stringstyle=\color{codestr},
  showstringspaces=false,
  frame=single,
  breaklines=true,
  tabsize=2,
  inputencoding=utf8,
  extendedchars=true
}
\lstset{style=pystyle}

\lstset{
    literate=
     {á}{{\'a}}1 {à}{{\`a}}1 {ã}{{\~a}}1 {â}{{\^a}}1
     {é}{{\'e}}1 {ê}{{\^e}}1
     {í}{{\'i}}1
     {ó}{{\'o}}1 {õ}{{\~o}}1 {ô}{{\^o}}1
     {ú}{{\'u}}1
     {ç}{{\c{c}}}1
     {Á}{{\'A}}1 {Ã}{{\~A}}1
     {É}{{\'E}}1 {Ê}{{\^E}}1
     {Í}{{\'I}}1
     {Ó}{{\'O}}1 {Õ}{{\~O}}1 {Ô}{{\^O}}1
     {Ú}{{\'U}}1
     {Ç}{{\c{C}}}1,
}


% =====================
% Título
% =====================
\title{\textbf{TP — Parte II: Aprendizagem Contínua via Otimização Min--Max}\\\large Disciplina de Otimização Não Linear}
\author{\textbf{Entrega}: Código (repositório), Relatório ($\leq$ 5 págs.), Apresentação (5 min)}
\date{\today}

\begin{document}
\maketitle

\section*{Motivação}
Treinar um modelo sequencialmente em tarefas distintas (A depois B) leva a \emph{esquecimento catastrófico}. A tensão entre \textbf{estabilidade} (reter A) e \textbf{plasticidade} (aprender B) pode ser modelada como um \textbf{problema min--max}. Neste trabalho, você implementa essa formulação com \textbf{passos alternados} e métodos obrigatórios de estabilização: \emph{OGDA} (Optimistic Gradient Descent--Ascent) e \emph{Extragradient}.

Usaremos \textbf{Split--MNIST}: Tarefa A = dígitos \{0--4\}, Tarefa B = \{5--9\}. O modelo base fornecido é uma CNN pequena.

\section*{Formulação de Otimização}
Denote por $\bm{\theta}$ os parâmetros e por $\bm{\theta}_A$ a solução após treinar a Tarefa A. Durante o aprendizado da Tarefa B, limitamos o esquecimento com a medida diferenciável de \textbf{deriva de parâmetros}:
\begin{equation}
F_{\text{drift}}(\bm{\theta};\bm{\theta}_A)=\tfrac{1}{2}\,\lVert \bm{\theta}-\bm{\theta}_A\rVert_2^2.
\end{equation}
Seja $\mathcal{L}_B(\bm{\theta})$ a perda em B. O jogo min--max é:
\begin{equation}
\min_{\bm{\theta}}\;\max_{\lambda\ge 0}\;\; \mathcal{L}_B(\bm{\theta})-\lambda\,F_{\text{drift}}(\bm{\theta};\bm{\theta}_A).
\label{eq:minmax}
\end{equation}
O minimizador busca baixo erro em B sem se afastar demais de $\bm{\theta}_A$; o maximizador $\lambda$ amplifica a penalização quando a deriva cresce. Após cada passo, projete $\lambda\leftarrow\max(0,\lambda)$.

\paragraph{OGDA (Obrigatório).}
Com $f(\bm{\theta},\lambda)=\mathcal{L}_B(\bm{\theta})-\lambda F_{\text{drift}}(\bm{\theta})$, o OGDA usa passo \emph{otimista} com gradiente corrente e anterior:
\begin{align}
\bm{\theta}_{t+1} &= \bm{\theta}_t - \eta_\theta\bigl(2\nabla_{\bm{\theta}} f_t - \nabla_{\bm{\theta}} f_{t-1}\bigr),\\
\lambda_{t+1} &= \bigl[\lambda_t + \eta_\lambda\bigl(2\nabla_{\lambda} f_t - \nabla_{\lambda} f_{t-1}\bigr)\bigr]_+.
\end{align}

\paragraph{Extragradient (Obrigatório).}
Método de duas etapas: (i) extrapolação com gradiente em $(\bm{\theta}_t,\lambda_t)$; (ii) correção usando gradiente no ponto extrapolado:
\begin{align}
(\hat{\bm{\theta}},\hat{\lambda}) &= (\bm{\theta}_t,\lambda_t) - (\eta_\theta,-\eta_\lambda)\odot \nabla f(\bm{\theta}_t,\lambda_t),\\
(\bm{\theta}_{t+1},\lambda_{t+1}) &= (\bm{\theta}_t,\lambda_t) - (\eta_\theta,-\eta_\lambda)\odot \nabla f(\hat{\bm{\theta}},\hat{\lambda}),\quad \lambda_{t+1}\leftarrow[\lambda_{t+1}]_+.
\end{align}

\section*{Objetivos de Aprendizagem}
\begin{itemize}
  \item Formular aprendizagem contínua como min--max.
  \item Implementar OGDA e Extragradient com projeção de $\lambda$.
  \item Quantificar e visualizar a frente \emph{estabilidade--plasticidade}.
\end{itemize}

\section*{Tarefas}
\begin{enumerate}
  \item \textbf{Dados e Modelo}: dividir MNIST em A:\,\{0--4\} e B:\,\{5--9\}; usar a CNN pequena fornecida.
  \item \textbf{Treino em A}: treinar por épocas fixas ou até acurácia-alvo; salvar $\bm{\theta}_A$ e um pequeno $\mathcal{V}_A$ (validação de A) para monitorar.
  \item \textbf{Treino Min--Max em B}: implementar \eqref{eq:minmax} com OGDA \emph{e} Extragradient, passos alternados e projeção de $\lambda$.
  \item \textbf{Métricas e Gráficos}: acurácia em A e B ao longo das iterações; $\lVert\bm{\theta}-\bm{\theta}_A\rVert$; trajetória de $\lambda$; curva estabilidade--plasticidade (acurácia final em B vs. esquecimento).
  \item \textbf{Ablações}: comparar OGDA vs. Extragradient; variar $\eta_\theta$, $\eta_\lambda$ e $\lambda_0$.
  \item \textbf{Baselines}: \emph{fine-tuning} ingênuo em B; regularização L2 com $\lambda$ fixo (sem jogo adversarial).
\end{enumerate}

\section*{Entrega e Avaliação (100 pts)}
\textbf{Código (45)}: implementação limpa PyTorch com flags; OGDA e Extragradient; semente determinística; plots salvos; README com comandos.\\
\textbf{Relatório (45)}: formulação, métodos, experimentos, plots, análise da frente e comportamento OGDA vs. Extragradient.\\
\textbf{Slides (10)}: apresentação de 5 min com figuras e conclusões.

\section*{Ética}
Usar apenas os dados e modelos fornecidos para a disciplina.

\section*{Como rodar}
\begin{verbatim}
python baseline_minmax.py --ogda --epochsA 3 --epochsB 3 --theta_lr 1e-3 --lambda_lr 1e-2
python baseline_minmax.py --extragrad --epochsA 3 --epochsB 3 --theta_lr 1e-3 --lambda_lr 1e-2
python run_experiments.py  # varredura simples
\end{verbatim}

\section*{Resumo Teórico (1 página)}
\textbf{Por que min--max?} Diferente de regularização com $\lambda$ fixo, o jogo permite que $\lambda$ se adapte via ascensão em resposta à deriva real. \textbf{OGDA} reduz ciclos usando passo otimista com gradientes corrente e anterior. \textbf{Extragradient} faz extrapolação e correção, estabilizando dinâmicas rotacionais típicas de jogos.\\
\textbf{O que plotar?} Acurácias em A/B, $\lVert\bm{\theta}-\bm{\theta}_A\rVert$, trajetória de $\lambda$ e a frente estabilidade--plasticidade ao variar hiperparâmetros.\\
\textbf{Mensagem-chave:} min--max transforma estabilidade vs. plasticidade em jogo de otimização; OGDA e Extragradient trazem estabilidade prática.

% =====================
% APÊNDICES COM CÓDIGO
% =====================
\clearpage
\appendix

\section*{Apêndice A: baseline\_minmax.py}
\begin{lstlisting}
# baseline_minmax.py
import math, os, random, argparse
from dataclasses import dataclass
import numpy as np
import torch
import torch.nn as nn
import torch.nn.functional as F
import torch.optim as optim
from torch.utils.data import DataLoader, Subset
from torchvision import datasets, transforms
import matplotlib.pyplot as plt

# -----------------------------
# Reproducibilidade
# -----------------------------
def set_seed(seed: int = 42):
    random.seed(seed)
    np.random.seed(seed)
    torch.manual_seed(seed)
    torch.cuda.manual_seed_all(seed)
    torch.backends.cudnn.deterministic = True
    torch.backends.cudnn.benchmark = False

# -----------------------------
# Dados
# -----------------------------
class SplitMNIST(torch.utils.data.Dataset):
    def __init__(self, root, train, digits, transform=None, download=True):
        self.base = datasets.MNIST(root=root, train=train,
                                   transform=transform, download=download)
        self.idx = [i for i, (_, y) in enumerate(self.base) if y in digits]
        self.digits = set(digits)
    def __len__(self):
        return len(self.idx)
    def __getitem__(self, i):
        x, y = self.base[self.idx[i]]
        return x, y

# -----------------------------
# Modelo
# -----------------------------
class SmallCNN(nn.Module):
    def __init__(self, num_classes=10):
        super().__init__()
        self.conv1 = nn.Conv2d(1, 16, 3, padding=1)
        self.conv2 = nn.Conv2d(16, 32, 3, padding=1)
        self.pool = nn.MaxPool2d(2, 2)
        self.fc1 = nn.Linear(32 * 7 * 7, 128)
        self.fc2 = nn.Linear(128, num_classes)
    def forward(self, x):
        x = self.pool(F.relu(self.conv1(x)))  # 28->14
        x = self.pool(F.relu(self.conv2(x)))  # 14->7
        x = x.view(x.size(0), -1)
        x = F.relu(self.fc1(x))
        return self.fc2(x)

# -----------------------------
# Utilitários
# -----------------------------
@torch.no_grad()
def evaluate(model, loader, device):
    model.eval()
    correct, total, loss_sum = 0, 0, 0.0
    ce = nn.CrossEntropyLoss()
    for x, y in loader:
        x, y = x.to(device), y.to(device)
        logits = model(x)
        loss = ce(logits, y)
        loss_sum += loss.item() * x.size(0)
        pred = logits.argmax(1)
        correct += (pred == y).sum().item()
        total += x.size(0)
    return loss_sum / total, correct / total

# Flatten parâmetros
def flatten_params(model):
    return torch.cat([p.detach().view(-1) for p in model.parameters()])

def flatten_params_grad(model):
    return torch.cat([p.grad.view(-1) if p.grad is not None else
                      torch.zeros_like(p).view(-1) for p in model.parameters()])

# Medida de deriva: 0.5 * ||theta - theta_A||^2
@torch.no_grad()
def param_drift(model, theta_A):
    theta = flatten_params(model)
    return 0.5 * torch.sum((theta - theta_A) ** 2).item()

# -----------------------------
# Treino na Tarefa A
# -----------------------------
def train_task_A(model, loader, device, epochs=3, lr=1e-3):
    opt = optim.Adam(model.parameters(), lr=lr)
    ce = nn.CrossEntropyLoss()
    for ep in range(epochs):
        model.train()
        for x, y in loader:
            x, y = x.to(device), y.to(device)
            opt.zero_grad()
            loss = ce(model(x), y)
            loss.backward()
            opt.step()
    theta_A = flatten_params(model).clone().detach()
    return theta_A

# -----------------------------
# Min  Max na Tarefa B
# -----------------------------
@dataclass
class MinMaxConfig:
    theta_lr: float = 1e-3
    lambda_lr: float = 1e-2
    epochs: int = 3
    ogda: bool = True  # se False, usa Extragradient
    batch_size: int = 128
    lambda_init: float = 0.0
    log_every: int = 100

class LambdaParam(torch.nn.Module):
    def __init__(self, init=0.0, device='cpu'):
        super().__init__()
        self.raw = nn.Parameter(torch.tensor([init], dtype=torch.float32,
                                             device=device))
    def value(self):
        return torch.clamp(self.raw, min=0.0)

# OGDA
def minmax_train_task_B_ogda(model, loader_B, loader_val_A, device, theta_A, cfg: MinMaxConfig):
    ce = nn.CrossEntropyLoss()
    lam = LambdaParam(cfg.lambda_init, device).to(device)
    opt_theta = optim.Adam(model.parameters(), lr=cfg.theta_lr)
    opt_lambda = optim.SGD(lam.parameters(), lr=cfg.lambda_lr)

    g_theta_prev = [torch.zeros_like(p) for p in model.parameters()]
    g_lambda_prev = torch.zeros(1, device=device)

    logs = []
    it = 0
    for ep in range(cfg.epochs):
        for xB, yB in loader_B:
            model.train()
            xB, yB = xB.to(device), yB.to(device)

            opt_theta.zero_grad(set_to_none=True)
            opt_lambda.zero_grad(set_to_none=True)

            logits_B = model(xB)
            L_B = ce(logits_B, yB)
            theta_vec = flatten_params(model)
            F_drift = 0.5 * torch.sum((theta_vec - theta_A) ** 2)
            obj = L_B - lam.value() * F_drift

            obj.backward()
            g_theta = [p.grad.detach().clone() if p.grad is not None
                       else torch.zeros_like(p) for p in model.parameters()]
            g_lambda = -F_drift.detach().clone().unsqueeze(0)  # grad ascendente

            with torch.no_grad():
                for p, gt, gtp in zip(model.parameters(), g_theta, g_theta_prev):
                    p.add_( -cfg.theta_lr * (2*gt - gtp) )
                lam.raw.add_( cfg.lambda_lr * (2*g_lambda - g_lambda_prev) )
                lam.raw.data.clamp_(min=0.0)

            g_theta_prev = [gt.detach().clone() for gt in g_theta]
            g_lambda_prev = g_lambda.detach().clone()

            if it % cfg.log_every == 0:
                model.eval()
                lossA, accA = evaluate(model, loader_val_A, device)
                logs.append({
                    'iter': it, 'ep': ep,
                    'L_B': L_B.item(), 'F_drift': F_drift.item(),
                    'lambda': lam.value().item(),
                    'valA_loss': lossA, 'valA_acc': accA,
                })
            it += 1
    return logs

# Extragradient
def minmax_train_task_B_extragradient(model, loader_B, loader_val_A, device, theta_A, cfg: MinMaxConfig):
    ce = nn.CrossEntropyLoss()
    lam = LambdaParam(cfg.lambda_init, device).to(device)
    theta_lr = cfg.theta_lr
    lambda_lr = cfg.lambda_lr
    logs = []
    it = 0

    for ep in range(cfg.epochs):
        for xB, yB in loader_B:
            model.train()
            xB, yB = xB.to(device), yB.to(device)

            # 1) Gradiente no ponto atual
            logits_B = model(xB)
            L_B = ce(logits_B, yB)
            theta_vec = flatten_params(model)
            F_drift = 0.5 * torch.sum((theta_vec - theta_A) ** 2)
            obj = L_B - lam.value() * F_drift

            g_theta = torch.autograd.grad(obj, list(model.parameters()),
                                          retain_graph=True, allow_unused=True)
            g_lambda = torch.autograd.grad(obj, lam.raw, retain_graph=True)[0]

            # 2) Extrapolação
            with torch.no_grad():
                for p, gt in zip(model.parameters(), g_theta):
                    if gt is not None:
                        p.copy_(p - theta_lr * gt)
                lam.raw.add_(lambda_lr * g_lambda)
                lam.raw.data.clamp_(min=0.0)

            # 3) Gradiente no ponto extrapolado
            logits_B2 = model(xB)
            L_B2 = ce(logits_B2, yB)
            theta_vec2 = flatten_params(model)
            F_drift2 = 0.5 * torch.sum((theta_vec2 - theta_A) ** 2)
            obj2 = L_B2 - lam.value() * F_drift2

            # 4) Correção
            g_theta2 = torch.autograd.grad(obj2, list(model.parameters()),
                                           retain_graph=False, allow_unused=True)
            g_lambda2 = torch.autograd.grad(obj2, lam.raw)[0]
            with torch.no_grad():
                for p, gt in zip(model.parameters(), g_theta2):
                    if gt is not None:
                        p.add_( -theta_lr * gt )
                lam.raw.add_( lambda_lr * g_lambda2 )
                lam.raw.data.clamp_(min=0.0)

            if it % cfg.log_every == 0:
                model.eval()
                lossA, accA = evaluate(model, loader_val_A, device)
                logs.append({
                    'iter': it, 'ep': ep,
                    'L_B': L_B2.item(), 'F_drift': F_drift2.item(),
                    'lambda': lam.value().item(),
                    'valA_loss': lossA, 'valA_acc': accA,
                })
            it += 1
    return logs

# -----------------------------
# Pipeline
# -----------------------------
def main():
    parser = argparse.ArgumentParser()
    parser.add_argument('--seed', type=int, default=42)
    parser.add_argument('--epochsA', type=int, default=3)
    parser.add_argument('--epochsB', type=int, default=3)
    parser.add_argument('--batch', type=int, default=128)
    parser.add_argument('--theta_lr', type=float, default=1e-3)
    parser.add_argument('--lambda_lr', type=float, default=1e-2)
    parser.add_argument('--ogda', action='store_true')
    parser.add_argument('--extragrad', action='store_true')
    parser.add_argument('--device', type=str,
                        default='cuda' if torch.cuda.is_available() else 'cpu')
    args = parser.parse_args()

    assert not (args.ogda and args.extragrad), 'Escolha um: --ogda ou --extragrad'
    if not args.ogda and not args.extragrad:
        args.ogda = True

    set_seed(args.seed)
    device = torch.device(args.device)

    transform = transforms.Compose([transforms.ToTensor()])
    trainA = SplitMNIST(root='./data', train=True, digits={0,1,2,3,4}, transform=transform)
    trainB = SplitMNIST(root='./data', train=True, digits={5,6,7,8,9}, transform=transform)
    valA_full = datasets.MNIST(root='./data', train=False, transform=transform, download=True)
    valA_idx = [i for i, (_, y) in enumerate(valA_full) if y in {0,1,2,3,4}]
    valA = Subset(valA_full, valA_idx)

    loaderA = DataLoader(trainA, batch_size=args.batch, shuffle=True, num_workers=2)
    loaderB = DataLoader(trainB, batch_size=args.batch, shuffle=True, num_workers=2)
    loaderValA = DataLoader(valA, batch_size=256, shuffle=False, num_workers=2)

    model = SmallCNN().to(device)

    print('Treinando Tarefa A...')
    theta_A = train_task_A(model, loaderA, device, epochs=args.epochsA, lr=1e-3)
    lossA, accA = evaluate(model, loaderValA, device)
    print(f'Após A: valA loss {lossA:.4f}, acc {accA:.4f}')

    cfg = MinMaxConfig(theta_lr=args.theta_lr, lambda_lr=args.lambda_lr,
                       epochs=args.epochsB, ogda=args.ogda, batch_size=args.batch)

    print('Treinando Tarefa B (min--max)...')
    if args.ogda:
        logs = minmax_train_task_B_ogda(model, loaderB, loaderValA, device, theta_A, cfg)
    else:
        logs = minmax_train_task_B_extragradient(model, loaderB, loaderValA, device, theta_A, cfg)

    os.makedirs('outputs', exist_ok=True)
    iters = [d['iter'] for d in logs]
    for key in ['L_B', 'F_drift', 'lambda', 'valA_loss', 'valA_acc']:
        plt.figure()
        plt.plot(iters, [d[key] for d in logs])
        plt.xlabel('iteration'); plt.ylabel(key); plt.title(key)
        plt.tight_layout()
        plt.savefig(os.path.join('outputs', f'{key}.png'))
        plt.close()

    print('Concluído. Plots em outputs/.')
if __name__ == '__main__':
    main()
\end{lstlisting}

\section*{Apêndice B: run\_experiments.py}
\begin{lstlisting}
# run_experiments.py
import os, subprocess

GRID = {
    'ogda': [True, False],      # True -> OGDA, False -> Extragradient
    'theta_lr': [5e-4, 1e-3, 2e-3],
    'lambda_lr': [5e-3, 1e-2, 2e-2],
}
EPOCHS_A = 3
EPOCHS_B = 3

os.makedirs('runs', exist_ok=True)

for ogda in GRID['ogda']:
    for tlr in GRID['theta_lr']:
        for llr in GRID['lambda_lr']:
            tag = f"{'ogda' if ogda else 'extragrad'}_tlr{tlr}_llr{llr}"
            cmd = [
                'python', 'baseline_minmax.py',
                '--epochsA', str(EPOCHS_A),
                '--epochsB', str(EPOCHS_B),
                '--theta_lr', str(tlr),
                '--lambda_lr', str(llr),
                '--batch', '128',
            ]
            if ogda:
                cmd.append('--ogda')
            else:
                cmd.append('--extragrad')
            print('Running', tag)
            subprocess.run(cmd, check=True)
\end{lstlisting}

\end{document}
